\begin{mistake}{2026-07-15}{09}

\begin{problemcontent}
已知曲线 \(y=y(x)\) 上任一点 \((x,y)\) 处的切线斜率为
\[
\frac{1}{x\sqrt{x^2-1}},
\]
且曲线通过点 \((-2,0)\)，求该曲线方程。
\end{problemcontent}

\begin{solutioncontent}
由题意，切线斜率即导数，故
\[
y'=\frac{dy}{dx}=\frac{1}{x\sqrt{x^2-1}}.
\]
两边积分：
\[
y=\int \frac{1}{x\sqrt{x^2-1}}\,dx.
\]
因曲线经过 \((-2,0)\)，应取定义分支 \(x<-1\)。计算
\[
\frac{d}{dx}\arcsin\frac1x
=\frac{\left(\frac1x\right)'}{\sqrt{1-\left(\frac1x\right)^2}}
=\frac{-\frac1{x^2}}{\sqrt{1-\left(\frac1x\right)^2}}.
\]
又
\[
\sqrt{1-\frac1{x^2}}
=\sqrt{\frac{x^2-1}{x^2}}
=\frac{\sqrt{x^2-1}}{|x|}.
\]
当 \(x<-1\) 时，\(|x|=-x\)，所以
\[
\frac{d}{dx}\arcsin\frac1x
=\frac{-\frac1{x^2}}{\frac{\sqrt{x^2-1}}{-x}}
=\frac{1}{x\sqrt{x^2-1}}.
\]
因此
\[
y=\arcsin\frac1x+C.
\]
代入 \((-2,0)\)：
\[
0=\arcsin\left(-\frac12\right)+C=-\frac{\pi}{6}+C,
\]
得
\[
C=\frac{\pi}{6}.
\]
故曲线方程为
\[
y=\arcsin\frac1x+\frac{\pi}{6},\quad x<-1.
\]
等价地，也可写为
\[
y=-\arccos\frac1x+\frac{2\pi}{3},\quad x<-1.
\]
\end{solutioncontent}

\mistakeknowledge{
\begin{itemize}
  \item \textbf{核心公式/定理}：切线斜率等于导数，即 \(y'=\frac{dy}{dx}\)；反三角函数求导需结合复合函数求导。
  \item \textbf{易错点}：\(\sqrt{x^2}=|x|\)，因此 \(\sqrt{1-\frac1{x^2}}=\frac{\sqrt{x^2-1}}{|x|}\)，不能直接把 \(|x|\) 写成 \(x\)。
  \item \textbf{关联知识}：在 \(x<-1\) 上，\(\frac{d}{dx}\arcsin\frac1x=\frac{1}{x\sqrt{x^2-1}}\)，而 \(\frac{d}{dx}\arccos\frac1x=-\frac{1}{x\sqrt{x^2-1}}\)。
\end{itemize}}

\end{mistake}
